% Methods: resolving the front/back ambiguity of body orientation (heading_correction_viterbi in utils/utils.py)
% Requires amsmath.

\subsection*{Resolving the front/back ambiguity of body orientation}

\paragraph{The problem.}
Body orientation was measured by fitting an ellipse to the fly's silhouette in a top-down camera. An ellipse
defines the body axis but not which end is the head, so in every frame $t = 1,\dots,n$ of a trajectory the
heading is either the measured axis $a_t$ or the opposite direction $a_t + \pi$. We write this binary choice as
$s_t \in \{0, 1\}$, with candidate heading
\begin{equation}
  \theta_t(s_t) = a_t + \pi s_t .
\end{equation}
Two kinds of information constrain the choice. First, physical cues point roughly forward: flies generate
aerodynamic force mostly along the body axis and usually fly forward through the air. Second, body orientation
changes little between consecutive frames (10\,ms apart), so neighbouring frames should usually make the same
choice. Neither is sufficient on its own. The physical cues are unreliable whenever the force or the airspeed is
small, and continuity only fixes the orientation of each frame \emph{relative} to its neighbours: following it
from a single starting decision lets one mistake propagate through the rest of the trajectory. We therefore
combine both kinds of information in a single score and choose the orientation of all frames of a trajectory
jointly.

\paragraph{Preparing the body axis.}
Because the axis is only defined modulo $\pi$, we first make it continuous in time by adding multiples of $\pi$
so that it never changes by more than $90^\circ$ between frames (\texttt{numpy.unwrap} with period $\pi$). A frame
whose axis differs by more than $40^\circ$ from both neighbours, while the neighbours lie within $30^\circ$ of each
other, is treated as a single bad ellipse fit and its axis is replaced by linear interpolation.

\paragraph{Physical cues and their reliability.}
The aerodynamic force (``thrust'') follows from the horizontal force balance $m\,\mathbf{a}_t = \mathbf{T}_t -
b\,\mathbf{v}_{\mathrm{air},t}$, with linear drag:
\begin{equation}
  \mathbf{T}_t = m\,\mathbf{a}_t + b\,\mathbf{v}_{\mathrm{air},t}, \qquad
  \mathbf{v}_{\mathrm{air},t} = \mathbf{v}_t - \mathbf{w},
\end{equation}
where $\mathbf{v}_t$ is the ground velocity, $\mathbf{w}$ the wind velocity, $\mathbf{a}_t$ the acceleration
(Savitzky--Golay derivative of $\mathbf{v}_t$), $m = 1.1$\,mg the body mass and $b = 11\,\mu\mathrm{N\,s\,m^{-1}}$
the drag coefficient of female \emph{D.~melanogaster} \citep{fuller2014}. We denote the directions of
$\mathbf{T}_t$ and $\mathbf{v}_{\mathrm{air},t}$ by $\varphi_{T,t}$ and $\varphi_{A,t}$.

The direction of a vector is dominated by measurement noise when its magnitude is small, so each cue must be
weighted by how informative it is. We measured this directly, using the body axis itself, which is reliable even
though its sign is not: in frames with a clearly elongated ellipse, the median angle between the axis and the
thrust direction (ignoring sign) was $38$--$43^\circ$ for thrust below $3\,\mu$N, close to the $45^\circ$ expected
if the two were unrelated, and $6$--$7^\circ$ above $8\,\mu$N. The airspeed direction behaved the same way:
$37$--$43^\circ$ below $0.3$\,m\,s$^{-1}$ and $2$--$3^\circ$ above $0.7$\,m\,s$^{-1}$. Each cue therefore receives a
weight that ramps linearly between these limits,
\begin{equation}
  r(x;\, x_{\mathrm{lo}}, x_{\mathrm{hi}}) = f + (1 - f)\,
  \min\!\Big(1, \max\!\Big(0, \frac{x - x_{\mathrm{lo}}}{x_{\mathrm{hi}} - x_{\mathrm{lo}}}\Big)\Big),
\end{equation}
with a small floor $f = 0.05$ so that weak cues can still break ties. The \emph{evidence} that frame $t$ provides for
candidate heading $\theta_t(s)$ is its agreement with the two cue directions:
\begin{equation}
  E_t(s) = w_{T,t}\cos\!\big(\theta_t(s) - \varphi_{T,t}\big)
         + \lambda\, w_{A,t}\cos\!\big(\theta_t(s) - \varphi_{A,t}\big),
\end{equation}
with weights $w_{T,t} = r(|\mathbf{T}_t|;\, 3, 8\,\mu\mathrm{N})$ and
$w_{A,t} = r(|\mathbf{v}_{\mathrm{air},t}|;\, 0.3, 0.7\,\mathrm{m\,s^{-1}})$. The cosine is $+1$ when a candidate points along the cue and $-1$ when it points against it, so the two candidates
of a frame receive opposite evidence, $E_t(1) = -E_t(0)$. The airspeed term is down-weighted ($\lambda = 0.3$)
because thrust already contains the air velocity through the drag term; it adds a slight bias toward forward
flight through the air. Ground-velocity direction is not used, because flies are often carried sideways or
backwards by the wind.

\paragraph{Continuity between frames.}
Consecutive frames should usually keep the same choice. We score this with
\begin{equation}
  C_t(s, s') = \kappa_t \cos\!\big(\theta_{t+1}(s') - \theta_t(s)\big)
             = \pm\, \kappa_t \cos(a_{t+1} - a_t),
\end{equation}
which is positive when frames $t$ and $t+1$ make the same choice ($s' = s$) and negative when they switch. A switch
therefore costs $2\kappa_t\cos(a_{t+1} - a_t)$: a lot when the axis barely moved, and nothing when the axis turned by
$90^\circ$ in one frame, where both interpretations of the turn are equally plausible. The strength $\kappa_t$
reflects how trustworthy the measured axis is. Nearly round ellipses (eccentricity, the minor-to-major axis ratio,
close to 1) give a poorly defined axis: changes of more than $45^\circ$ between consecutive frames occurred in
$0.2\%$ of frame pairs when both ellipses had eccentricity below $0.6$, but in $20\%$ when either exceeded $0.9$.
Hence
\begin{equation}
  \kappa_t = \frac{\kappa_{\min} + (\kappa_{\max} - \kappa_{\min})\,\big(1 - \rho_t\big)}{g_t}, \qquad
  \rho_t = \min\!\Big(1, \max\!\Big(0, \frac{\max(e_t, e_{t+1}) - 0.6}{0.3}\Big)\Big),
\end{equation}
with $\kappa_{\max} = 5$, $\kappa_{\min} = 0.2$, $e_t$ the eccentricity, and $g_t \ge 1$ the number of frame intervals
spanned by the step (larger than 1 only across missing frames).

\paragraph{Choosing all orientations jointly.}
The score of a complete sequence of choices $\mathbf{s} = (s_1, \dots, s_n)$ is the total evidence plus the total
continuity,
\begin{equation}
  S(\mathbf{s}) = \sum_{t=1}^{n} E_t(s_t) + \sum_{t=1}^{n-1} C_t(s_t, s_{t+1}),
\end{equation}
and the corrected heading is given by the sequence with the highest score, $\hat{\mathbf{s}} = \arg\max_{\mathbf{s}}
S(\mathbf{s})$. Maximizing over whole sequences is what makes the method robust. Strong evidence anywhere in a
trajectory is carried, through well-measured continuity, to frames whose own evidence is weak. Where the axis is
poorly measured, continuity is weak, so the parts of the trajectory on either side are each decided by their own
evidence instead of inheriting a single early decision. A frame is reversed relative to its neighbours only if
the evidence for doing so outweighs the cost of the two switches it requires.

A trajectory of $n$ frames has $2^n$ possible sequences, far too many to compare one by one. Because every term of
$S$ involves at most two consecutive frames, however, the best sequence can be found exactly with the Viterbi
algorithm \citep{viterbi1967}, a form of dynamic programming. Moving forward through the trajectory, we keep, for
each of the two choices at frame $t$, the score $V_t(s)$ of the best partial sequence that ends with that choice,
\begin{equation}
  V_1(s) = E_1(s), \qquad
  V_{t+1}(s') = E_{t+1}(s') + \max_{s} \big[ V_t(s) + C_t(s, s') \big],
\end{equation}
and remember which choice at frame $t$ achieved the maximum. Any best sequence that passes through choice $s'$ at
frame $t+1$ must reach it through the best partial sequence ending at frame $t$, so only these two running scores
are needed. The best final choice $\arg\max_s V_n(s)$ is then traced back through the remembered predecessors.
This yields the exact optimum in time proportional to $n$.

\paragraph{Confidence.}
To express how strongly the data support each frame's orientation, we interpret $\exp S(\mathbf{s})$ as the
relative plausibility of sequence $\mathbf{s}$, as in a hidden Markov model,
\begin{equation}
  P(\mathbf{s}) = \frac{\exp S(\mathbf{s})}{\sum_{\mathbf{s}'} \exp S(\mathbf{s}')},
\end{equation}
and define the confidence of frame $t$ as the total plausibility of all sequences that agree with the chosen
orientation at that frame,
\begin{equation}
  c_t = P(s_t = \hat s_t) = \frac{\sum_{\mathbf{s}:\, s_t = \hat s_t} \exp S(\mathbf{s})}
                                 {\sum_{\mathbf{s}} \exp S(\mathbf{s})}.
\end{equation}
These sums over all $2^n$ sequences are computed exactly with the forward--backward algorithm, which has the same
structure as the Viterbi recursion with the maximum replaced by a sum:
\begin{align}
  \alpha_1(s) &= e^{E_1(s)}, &
  \alpha_{t+1}(s') &= e^{E_{t+1}(s')} \sum_{s} \alpha_t(s)\, e^{C_t(s, s')}, \\
  \beta_n(s) &= 1, &
  \beta_t(s) &= \sum_{s'} e^{C_t(s, s') + E_{t+1}(s')}\, \beta_{t+1}(s'),
\end{align}
so that $c_t = \alpha_t(\hat s_t)\,\beta_t(\hat s_t) \big/ \sum_s \alpha_t(s)\,\beta_t(s)$. Here $\alpha_t$ summarizes the
evidence from the start of the trajectory up to frame $t$, and $\beta_t$ the evidence after it. A frame isolated from
its neighbours has $c_t = 1/(1 + e^{-2E_t(\hat s_t)})$, at most $0.88$ even for strong thrust pointing exactly along
the chosen heading; confidence close to 1 arises when well-measured continuity links a frame to many frames with
consistent evidence. Confidence is low where cues are weak or conflicting, or where nearly round ellipses cut a
frame off from strong evidence. Because multiplying all weights by a common factor leaves $\hat{\mathbf{s}}$
unchanged but changes $c_t$, the confidence ranks frames by reliability but is not a calibrated probability. In
rare frames it falls below $0.5$, where the best sequence as a whole disagrees with the most plausible choice
for that frame alone.

\paragraph{Parameters and validation.}
The parameters ($\lambda$, $\kappa_{\min}$, $\kappa_{\max}$, the ramp limits and the floor) were set from the
reliability analysis above, before any hand labels were collected. To test the method, one of us labelled
the front/back orientation of 285 trajectories (3{,}928 frames, about 14 per trajectory, spaced evenly along the
path) in a custom tool that showed the trajectory with the body ellipse, air velocity, thrust and a candidate
heading, without indicating which method had produced it. The sample over-represented trajectories on which the
Viterbi correction and a simpler procedure (aligning the first frames and then the whole trajectory with the
thrust direction, and otherwise following frame-to-frame continuity) disagreed, and
dataset-wide estimates therefore weight each sampling stratum by its size. The raw ellipse orientation had the
correct front/back direction in $79.8\%$ of frames (95\% bootstrap interval over trajectories, $74.9$--$84.8\%$),
the simpler procedure in $96.7\%$ ($96.1$--$97.3\%$) and the Viterbi correction in $97.9\%$ ($97.4$--$98.4\%$).
A search over 12{,}600 parameter combinations, evaluated by five-fold cross-validation over trajectories, did not
improve on the a priori parameters on held-out labels, so these were kept. Disagreement with the hand labels fell
monotonically with confidence, from $27\%$ of frames with $c_t < 0.7$ to $9\%$ for $0.9 \le c_t < 0.99$ and $1.5\%$ for
$c_t \ge 0.999$; $76\%$ of all frames had $c_t \ge 0.999$ and $10\%$ had $c_t < 0.9$. On ambiguous frames the
labeller tended to accept whichever orientation was shown. This pulls the measured accuracies of the two
corrections toward each other, so their difference is a conservative estimate, but it may flatter both relative
to the raw orientation, which was never shown.

% Bibliography entries assumed:
%   fuller2014: Fuller SB, Straw AD, Peek MY, Murray RM, Dickinson MH (2014) Flying Drosophila stabilize their
%               vision-based velocity controller by sensing wind with their antennae. PNAS 111:E1182-E1191.
%   viterbi1967: Viterbi AJ (1967) Error bounds for convolutional codes and an asymptotically optimum decoding
%               algorithm. IEEE Trans Inf Theory 13:260-269.
